\medskip

Dans l'espace muni d'un repère orthonormé d'unité $1 \mathrm{~cm}$, on considère les points

\begin{center}$\mathrm{D}(3~;~1~;~5)$, \qquad 
$\mathrm{E}(3~;~-2~;~-1)$, \qquad 
$\mathrm{F}(-1~;~2~;~1)$, \qquad 
$\mathrm{G}(3~;~2~;~-3)$.\end{center}

\smallskip

\begin{enumerate}
\item 
	\begin{enumerate}
		\item Déterminer les coordonnées des vecteurs $\overrightarrow{\text{EF}}$ et $\overrightarrow{\text{FG}}$.
		\item Justifier que les points $\mathrm{E}, \mathrm{F}$ et $\mathrm{G}$ ne sont pas alignés.
	\end{enumerate}
\item 
	\begin{enumerate}
		\item Déterminer une représentation paramétrique de la droite (FG).
		\item On appelle $\mathrm{H}$ le point de coordonnées $(2~;~2~;~-2)$.

Vérifier que $\mathrm{H}$ est le projeté orthogonal de E sur la droite (FG) .
		\item Montrer que l'aire du triangle EFG est égale à $12\mathrm{~cm}^{2}$.
	\end{enumerate}
\item 
	\begin{enumerate}
		\item Démontrer que le vecteur $\vec{n}\begin{pmatrix}2 \\ 1 \\ 2\end{pmatrix}$ est un vecteur normal au plan (EFG).
		\item Déterminer une équation cartésienne du plan (EFG) .
		\item Déterminer une représentation paramétrique de la droite $(d)$ passant par le point $\mathrm{D}$ et orthogonale au plan (EFG) .
		\item On note K le projeté orthogonal du point D sur le plan (EFG).

À l'aide des questions précédentes, calculer les coordonnées du point $\mathrm{K}$.
	\end{enumerate}
\item 
	\begin{enumerate}
		\item Vérifier que la distance $D K$ est égale à $5 \mathrm{~cm}$.
		\item En déduire le volume du tétraèdre DEFG.
	\end{enumerate}
\end{enumerate}

\bigskip
